
\renewcommand{\arraystretch}{1.5} % Default value: 1
%%%%%%%%%%%%%%
% Introduction
\section{Introduction}
The Mach-Zehnder interferometer is a simple but powerful tool used to show how light waves interfere with each other. It works by splitting a single beam of light into two separate paths using a "beamsplitter". These two beams travel along different routes before being brought back together by a second beamsplitter. Depending on what happens to the light on those two paths, the beams will either add up (constructive interference) or cancel each other out (destructive interference) when they reach the detectors.
%%%%%%%%%%%%%%%%%%%%%%%%%%%%%%%%%%%%%
\section{Objective}
The objective of this experiment is to learn the alignment of the Mach-Zehnder interferometer and observe the fringes. Also, measuring refractive index of air by using fringe shifts in a Mach-Zehnder interferometer is another objective of this experiment.
%%%%%%%%%%%%%%%%%%%%%%%%%%%%%%%%%%%%%
\section{Equipment Required}
The following equipments are required to perform this experiment:
\begin{itemize}
    \item $650 \text{ nm}$ wavelength Laser source,
    \item Two 50-50 Beam splitters,
    \item Two reflectors (mirrors),
    \item Two lenses for focusing,
    \item Two white screens for marking and measuring the shift of the interference fringes,
    \item The optical bench, optical mounts and holders,
    \item A lens for producing circular fringes,
    \item A vacuum cell, and a vacuum gauge with a hand pump.
\end{itemize}
%%%%%%%%%%%%%%%%%%%%%%%%%%%%%%%%%%%%%
\section{Theory}
The Mach-Zehnder interferometer operates on the principle of division of amplitude. Unlike division of wavefront (like Young's double-slit), this method takes a single light beam and splits its energy (amplitude) into two paths using a semi-transparent mirror, known as a beamsplitter. These two beams travel along separate paths (the upper and lower arms) before being recombined by a second beamsplitter.

\begin{figure}[H]
    \centering
    \includegraphics[width = 0.8\textwidth]{Image/Mach_Zehnder interferometer_schematic.jpg}
    \caption{Mach-Zehnder interferometer schematic}
    \label{fig:MZ_sch}
\end{figure}

The intensity of the light at the output detectors depends on the phase difference ($\Delta\phi$) between the two beams when they meet. If the waves arrive in phase, they interfere constructively; if they arrive $\pi$ out of phase, they interfere destructively.
\subsection{Phase Shifts}
A critical part of understanding this interferometer is the phase change that occurs during reflection. Standard physics "lore" suggests that every reflection causes a $\pi$ phase shift, but this is an oversimplification that violates the law of conservation of energy in this specific setup.

To resolve this, we must look at the Fresnel Equations for a dielectric beamsplitter:
\begin{itemize}
    \item \emph{External Reflection (Air to Dielectric):} When light reflects off the front surface (moving from a low refractive index to a high refractive index), it undergoes a phase shift of $\pi$.
    \item \emph{Internal Reflection (Glass to Dielectric):} When light reflects from "behind" the coating while inside the glass substrate (moving from high to low index), there is no phase shift ($0$).
    \item \emph{Transmission:} Light passing through the beamsplitter picks up no phase shift other than the delay caused by the optical path length of the glass.
\end{itemize}

By arranging the second beamsplitter such that its reflecting surface is on the "back" side relative to one path, we ensure that one output path experiences one more $\pi$ shift than the other. This explains why the two outputs (A and B) are always $\pi$ out of phase: when Detector A shows a bright fringe (constructive), Detector B must show a dark fringe (destructive).

\subsection{Mathematical Analysis of Phase Difference}
To understand the device's operation, we track the phase along the Upper ($l_1$) and Lower ($l_2$) paths. Let $t$ represent the optical path length through the beamsplitter glass (see \cref{fig:MZ_sch}).

\subsection{Path to Detector A}
\begin{itemize}
    \item \textbf{Upper Path to A:} Reflection at BS1 (Air-to-Dielectric: $+\pi$), Reflection at $100\%$ Mirror ($+\pi$), Transmission through BS2 ($0$). Total: $2\pi + \frac{2\pi(l_1 + t)}{\lambda}$.
    \item \textbf{Lower Path to A:} Transmission through BS1 ($0$), Reflection at $100\%$ Mirror ($+\pi$), Reflection at BS2 (Internal: $+\pi$ per manual symmetry). Total: $2\pi + \frac{2\pi(l_2 + t)}{\lambda}$.
\end{itemize}
The phase difference at A is: $\Delta\Phi_A = \frac{2\pi(l_1 - l_2)}{\lambda} = \delta$.

\subsection{Path to Detector B}
\begin{itemize}
    \item \textbf{Upper Path to B:} Identical to A, with no extra phase due to reflection at BS2 (internal) ($0$). Total: $2\pi + \frac{2\pi(l_1 + 2t)}{\lambda}$.
    \item \textbf{Lower Path to B:} Similar to A and transmission through BS2 ($0$). Total: $\pi + \frac{2\pi(l_2 + 2t)}{\lambda}$.
\end{itemize}
The phase difference at B is: $\Delta\Phi_B = \pi + \frac{2\pi(l_1 - l_2)}{\lambda} = \pi + \delta$.

\subsection{Refractive Index and Pressure}

According to the Lorenz-Lorentz equation, the refractive index of a gas is related to its density, which is directly proportional to its pressure ($P$) at a constant temperature. For air at atmospheric pressures, the relationship is approximately linear:


$$n - 1 \propto P$$

By evacuating a vacuum cell of length $d$ and counting the number of fringes $\Delta m$ that pass a reference point, we can calculate the slope of the refractive index relative to pressure:


$$\frac{n_i - n_f}{P_i - P_f} = \frac{\Delta m \lambda_0 / d}{P_i - P_f}$$

Assuming the refractive index of a vacuum is exactly $1$, we can use this slope to determine the refractive index of air at $1$ atmosphere ($760 \text{ mmHg}$ or $101.3 \text{ kPa}$).
%%%%%%%%%%%%%%%%%%%%%%%%%%%%%%%%%%%%%
\section{Procedure}
\begin{figure}[H]
    \centering
    \includegraphics[width = 0.7\textwidth]{Image/MZ_1.JPG}
    \caption{Mach-Zehnder interferometer setup (BS stands for beam splitter)}
    \label{fig:MZ_lab_1}
\end{figure}

The crucial part of this experiment is to align the setup for the interferometer (see \cref{fig:MZ_lab_1}). We used the following procedure to align it. And later on we observed the fringes and do the subsequent procedure to complete the experiment.
\begin{enumerate}
    \item \textbf{Initial Alignment and Beam Path}
    \begin{itemize}
    \item \textbf{Laser Alignment:} We began by adjusting the height of the laser source to ensure the beam travelled perfectly parallel to the optical table.
    \item \textbf{Beamsplitter Placement:} We positioned the first beamsplitter (BS1) and checked the reflecting side to ensure the dielectric coating faced the incoming beam.
    \item \textbf{Table Orientation:} We aligned the component with the mounting holes on the optical table to maintain a precise $90^\circ$ geometry.
    \item \textbf{Height Verification:} Using a board at the far end of the table, we checked for any horizontal or vertical shifts in the beam to ensure it remained level over long distances.
\end{itemize}
    \item \textbf{Mirror Alignment (Near-Far Method)} To ensure the two arms of the interferometer were perfectly aligned, we used the "Near-Far" method for the mirrors:
    \begin{itemize}
        \item \textbf{Far-Point Setting:} We first adjusted the angular knobs of the mirror to fix the beam position at a "far" reference point.
        \item \textbf{Near-Point Correction:} We then checked the "near" position; if it was not centred, we shifted the mirror mount spatially by hand rather than using the knobs.
        \item \textbf{Iteration:} We iteratively checked the far point and adjusted the angular position, then checked the near point and adjusted the spatial position, until the beam was perfect at both distances.
    \end{itemize}
    \item \textbf{Recombination and Fringe Setup}
    \begin{itemize}
        \item \textbf{Second Beamsplitter:} We placed the second beamsplitter at the intersection of the two paths and secured it firmly with mounting screws.
        \item \textbf{Beam Expansion:} We inserted a lens in front of the laser to expand the beam, allowing us to see the interference pattern more clearly.
        \item \textbf{Circular fringes}: We found circular fringes by putting a lens in one of the arm, also verified that the fringes at output A and output B were out of phase by $\pi$, ensuring that a bright center at one detector corresponded to a dark center at the other.
    \end{itemize}
    \item \textbf{Refractive Index Measurement}
    \begin{itemize}
        \item \textbf{Cell Integration:} We inserted a pressure cell into one arm of the interferometer (see \cref{fig:MZ_lab_2})
        \item \textbf{Pressure Variation:} We used a hand-pump to evacuate the cell while carefully recording the readings from the vacuum gauge.
        \item \textbf{Fringe Counting:} We marked a reference point on the screen and counted the number of fringes  that passed that point as the air was evacuated.
    \end{itemize}
\end{enumerate}

\begin{figure}[H]
    \centering
    \includegraphics[width = 0.7\textwidth]{Image/MZ_2.JPG}
    \caption{Mach-Zehnder interferometer setup with pressure cell (BS stands for beam splitter)}
    \label{fig:MZ_lab_2}
\end{figure}
%%%%%%%%%%%%%%%%%%%%%%%%%%%%%%%%%%%%%

\section{Data} 
Now after setting up the interferometer we got fringe patterns which are given bellow.

\begin{figure}[H]
     \centering
     % First Subfigure
     \begin{subfigure}[b]{0.45\textwidth}
         \centering
         \includegraphics[height=5cm, width=\textwidth]{Image/MZ_fringe_1.jpeg}
         \caption{Fringe from the left lens in \cref{fig:MZ_lab_1}}
         \label{fig:left_image_1}
     \end{subfigure}
     \hfill 
     % Second Subfigure
     \begin{subfigure}[b]{0.45\textwidth}
         \centering
         \includegraphics[height=5cm, width=\textwidth]{Image/MZ_fringe_2.jpeg}
         \caption{Fringe from the top lens in \cref{fig:MZ_lab_1}}
         \label{fig:right_image_1}
     \end{subfigure}
     
     \caption{Fringes without lens in one arm}
     \label{fig:finges_without_L}
\end{figure}
Now we put a lens in one of the arms and we got circular fringes.

\begin{figure}[H]
     \centering
     % First Subfigure
     \begin{subfigure}[b]{0.45\textwidth}
         \centering
         \includegraphics[height=5cm, width=\textwidth]{Image/MZ_fringe_circular_2.jpeg}
         \caption{Fringe from left lens}
         \label{fig:left_image_2}
     \end{subfigure}
     \hfill 
     % Second Subfigure
     \begin{subfigure}[b]{0.45\textwidth}
         \centering
         \includegraphics[height=5cm, width=\textwidth]{Image/MZ_fringe_circular_1.jpeg}
         \caption{Fringe from top lens}
         \label{fig:right_image_2}
     \end{subfigure}
     
     \caption{Fringes with lens in one arm}
     \label{fig:fringes_with_L}
\end{figure}

After that, in \emph{Part B} we put the pressure cell (see \cref{fig:MZ_lab_2}) and take three set of data of fringe shifts to plot the graph between number of fringe changes and pressure changes.

\subsubsection*{Set I}
\begin{table}[H]
\centering
\begin{tabular}{|c|c|}
\toprule
Pressure (mm Hg) & No of fringes shifted \\ \midrule
6               & 0                     \\ \hline
40              & 1                     \\ \hline
66              & 2                     \\ \hline
88              & 3                     \\ \hline
120             & 4                     \\ \hline
144             & 5                     \\ \hline
170             & 6                     \\ \hline
190             & 7                     \\ \hline
220             & 8                     \\ \hline
248             & 9                     \\ \hline
276             & 10                   \\
\bottomrule
\end{tabular}
\end{table}


\begin{figure}[H]
    \begin{center}
        \includegraphics[width = 0.8\textwidth]{../data analysis/Plot1.pdf}
    \end{center}
\end{figure}

\subsubsection*{Set II}

\begin{table}[H]
\centering
\begin{tabular}{|c|c|}
\toprule
Pressure (mm Hg) & No of fringes shifted \\\midrule
6               & 0                     \\\hline
36              & 1                     \\\hline
60              & 2                     \\\hline
88              & 3                     \\\hline
116             & 4                     \\\hline
144             & 5                     \\\hline
162             & 6                     \\\hline
190             & 7                     \\\hline
214             & 8                     \\\hline
242             & 9                     \\\hline
270             & 10                    \\
\bottomrule
\end{tabular}
\end{table}

\begin{figure}[H]
    \begin{center}
        \includegraphics[width = 0.8\textwidth]{../data analysis/Plot2.pdf}
    \end{center}
\end{figure}

\subsubsection*{Set III}

\begin{table}[H]
\centering
\begin{tabular}{|c|c|}
\toprule
Pressure (mm Hg) & No of fringes shifted \\\midrule
6               & 0                     \\\hline
30              & 1                     \\\hline
58              & 2                     \\\hline
80              & 3                     \\\hline
106             & 4                     \\\hline
134             & 5                     \\\hline
152             & 6                     \\\hline
180             & 7                     \\\hline
210             & 8                     \\\hline
232             & 9                     \\\hline
260             & 10                    \\
\bottomrule
\end{tabular}
\end{table}

\begin{figure}[H]
    \begin{center}
        \includegraphics[width = 0.8\textwidth]{../data analysis/Plot3.pdf}
    \end{center}
\end{figure}
%%%%%%%%%%%%%%%%%%%%%%%%%%%%%%%%%%%%%
\section{Analysis}
The list of slope from the above graphs: 
\begin{enumerate}
    \item \textbf{SET I}: The slope is $0.0379 \pm 0.0004$
    \item \textbf{SET II}: The slope is $0.0385 \pm 0.0004$
    \item \textbf{SET III}: The slope is $0.0395 \pm 0.0004$
\end{enumerate}
Based on the data we have in the slope, the average slope is calculated. 
\[
\text{Avg slope } = \frac{0.0379+0.385+0.0395}{3} = 0.0386
\]
Now using the formula mentioned in the theory part we have to calculate the refractive index. 
\[
\frac{\Delta n}{\Delta P} = \frac{n_f - n_i}{P_f - P_i} = \frac{\Delta m}{\Delta P} \times \frac{\lambda_i}{d} = \text{slope} \times \frac{\lambda_i}{d}
\]
Now here, $\lambda = 655 \; \text{nm}$ and $d = 7.9 \; \text{cm}$. Putting this into main equation we have, 

\[
\frac{n_f - n_i}{P_f - P_i} = 0.0386 \; \text{mm Hg}^{-1}\times \frac{655 \times 10^{-7} \; \text{cm}}{7.9 \; \text{cm}} = 3.203 \times 10 ^{-7} \; \text{mm Hg}^{-1}
\]
Now we use the standard assumption that at $P_i = 0 atm$, $n_i = 1$ which is the refractive index at the vacuum. Now we have to find the refractive index $(n_f)$ at the pressure $P_f = 1 \; \text{atm} = 760 \; \text{mm Hg}$
\begin{align*}
    n_f &= n_i + P_f \times 3.203 \times 10 ^{-7} \; \text{mm Hg}^{-1} \\
    &= 1 + (760 \; \text{mm Hg} \times 3.203 \times 10 ^{-7} \; \text{mm Hg}^{-1}) \\
    & = 1 + 0.000243 = 1.000243
\end{align*}
So we find that the refractive index for air at $1 \; \text{atm}$ pressure is 1.000243 at room temperature.
%%%%%%%%%%%%%%%%%%%%%%%%%%%%%%%%%%%%%
\section{Error analysis}
\subsection{Maximum Propagation of Error}
We have the main formula 
\begin{align*}
    &\frac{\Delta n}{\Delta P} = \frac{n_f - n_i}{P_f - P_i} = \frac{\Delta m}{\Delta P} \times \frac{\lambda_i}{d} = \text{slope} \times \frac{\lambda_i}{d} \\
    \Rightarrow & \Delta n = \Delta P \times \text{slope} \times \frac{\lambda_0}{d}
\end{align*}
Now taking the natural log,
\[
\ln (\Delta n) = \ln (\Delta P) + \ln (\text{slope}) + \ln \lambda_0 - \ln d
\]
Taking the derivative
\[
\frac{\delta \Delta n}{\Delta n} = \frac{\delta \Delta P}{\Delta P} + \frac{\delta (\text{slope})}{\text{slope}} - \frac{\delta d}{d}
\]
Considering the maximum probable error
\[
\left|\frac{\delta \Delta n}{\Delta n}\right| = \left|\frac{\delta \Delta P}{\Delta P}\right| + \left|\frac{\delta (\text{slope})}{\text{slope}}\right| + \left|\frac{\delta d}{d}\right|
\]
Now, $\delta \Delta P = 0$ for our case and $\delta d = 0.1 \; \text{cm}$ and $\delta (\text{slope}) = 0.0004$. Using this values,
\[
\left|\frac{\delta \Delta n}{\Delta n}\right| = \left|\frac{0.0004}{0.0386}\right| + \left|\frac{0.1}{7.9}\right| = 0.023
\]
Now, using $\Delta n = n_f - n_i$ and $\delta \Delta n = \delta n_f - \delta n_i = \delta n_f$, (as we used a constant value for $n_i$)
\begin{align*}
    \left|\frac{\delta n_f}{n_f - n_i}\right| &= 0.023 \\
    \delta n_f = (1.000243 - 1.000000) &  \times 0.023 = 0.000006
\end{align*}

So the error is $\pm 0.000006$, hence our final value for the the refractive index for air at $1 \; \text{atm}$ pressure is $1.000243 \pm 0.000006$ at room temperature.

\subsection{Theoretical Error}
From the internet we got that \textbf{for $\mathbf{15^{\circ}C}$, n = 1.000277 for $\mathbf{1}$ atm air pressure}. So the percentage error from the theoretical error is
\[
\% \; \text{error} = \frac{(1.000277 - 1.000243)}{1.000277} \times 100 \% = 0.0033 \%
\]
\section{Source of Error}
There are few sources of errors are given bellow
\begin{enumerate}
    \item Change of pressure in the vapour cell is not very stable as we have to manually pump air into it, which causes error in taking reading. Because of this ,there is error in observing fringe shift and taking reading of the pressure in the dial.
    \item The beam-splitters may be not completely $50:50$. which affect the contrast of the fringes.
\end{enumerate}
\section{Discussion}
Here are few key things which can be discussed here.
\begin{itemize}
    \item \textbf{Why Fringes}: As I mentioned in the theory that for the perfect alignment we should get only constructive or destructive interference over the screen so only bright or only dark spot should be there. But we saw fringes. The reason is I think misalignment of the beam. It can be either horizontal or vertical misalignment. Based on this misalignment, we can have either vertical or horizontal fringes respectively. 
    \item \textbf{Fringes without merging of beam}: We observed that even if we blocked the one arm, some fringes are non removable for some setup, we assumed that, this is because of some kind or dirt placed on the mirror or the beam splitter which is causing circular type of fringe pattern.
    \item \textbf{Small setup easy to align}: We observed that the smaller setup are easy to align. The reason is for larger setup it is difficult to control the divergence of beam and we are not using collimated beam so little divergence angle can cause problem for larger setup.  
\end{itemize}
%%%%%%%%%%%%%%%%%%%%%%%%%%%%%%%%%%%%%
\section{Conclusion}
We successfully observed the linear as well as circular fringes using the Mach-Zehnder Interferometer. We found that in one arm the circular fringe have dark center and in another arm the circular fringe have bright centre. Also using the Pressure cell and fringe shift measurement we found the refractive index of the air at 1 atm pressure at room temperature. Hence we have, the refractive index for air at $1 \; \text{atm}$ pressure is $1.000243 \pm 0.000006$ at room temperature.

%%%%%%%%%%%%%%%%%%%%%%%%%%%%%%%%%%%%%